% GENERATED FILE. Edit canonical lessons, then run npm run generate:books.
% canonical-source-hash: fnv1a64:c8584d928121d787
% canonical-lessons: JA-C171-saisho, JA-C171-mazu, JA-C171-tsugi, JA-C171-saigo, JA-C171-sorosoro, JA-R171-review-170-171, JA-R171-review-final-171

\chapter{First of All, Next and Last}
\label{ch:ja-first-of-all-next-and-last}
\hlchaptermodality{171}

\begin{chapteropening}
\textbf{By the end of this chapter:} \emph{I can put my plans in order.}

Complete the last lesson of chapter 171: I can put my plans in order.
\end{chapteropening}

\section[\texorpdfstring{\ja{さいしょ} (saisho)}{さいしょ (saisho)}]{\ja{さいしょ} (saisho) — the beginning; at first}

\label{lesson:JA-C171-saisho}



\begin{quote}
Before the new one: say the Japanese for to celebrate, then the Japanese for to spend time.
\end{quote}

\subsection*{You'll want to know: \ja{さいしょ}}
\textbf{\ja{さいしょ}} — \emph{saisho} — \textquotedblleft{}the beginning; at first\textquotedblright{}.

\begin{quote}
\textbf{\ja{あしたは} \ja{さいしょに} \ja{えきに} \ja{いきます。}} — \emph{ashita wa saisho ni eki ni ikimasu} — \textquotedblleft{}Tomorrow I will go to the station first.\textquotedblright{}
\end{quote}

\begin{cousinweb}
\textbf{\ja{さいしょ}} (\emph{saisho}) is Sino-Japanese (kango): two Chinese characters read with their on'yomi, readings borrowed from Chinese. They are \emph{sai}, \textquotedblleft{}most, utmost\textquotedblright{} (Middle Chinese \emph{tswajH}) and \emph{sho}, \textquotedblleft{}first time, beginning\textquotedblright{} (\emph{tsrhjo}) (KANJIDIC2; Middle Chinese in Baxter's transcription).
\end{cousinweb}

\begin{grammarlens}[title={What you've built}]
One more word for everyday life at A2.
\end{grammarlens}

\subsection*{Your turn}
\begin{itemize}
\raggedright
  \item \emph{Say it:} \emph{saisho}
  \item \emph{Say it:} \emph{saisho}, once more
  \item \emph{Say it:} \emph{ashita wa saisho ni eki ni ikimasu} — Tomorrow I will go to the station first.
  \item \emph{Recall:} say the Japanese for to celebrate, then the Japanese for to spend time, then say \emph{saisho} again
\end{itemize}

\begin{tcolorbox}[breakable,colback=teal!4,colframe=teal!35!black,title={Before you move on}]
What does \ja{さいしょ} mean? (\textquotedblleft{}The beginning; at first\textquotedblright{}.) Say it once more.
\end{tcolorbox}
\section[\texorpdfstring{\ja{まず} (mazu)}{まず (mazu)}]{\ja{まず} (mazu) — first of all; to begin with}

\label{lesson:JA-C171-mazu}



\begin{quote}
Before the new one: say the Japanese for to spend time, then the Japanese for the beginning.
\end{quote}

\subsection*{You'll want to know: \ja{まず}}
\textbf{\ja{まず}} — \emph{mazu} — \textquotedblleft{}first of all; to begin with\textquotedblright{}.

\begin{quote}
\textbf{\ja{あしたは} \ja{まず} \ja{かおを} \ja{あらいます。}} — \emph{ashita wa mazu kao o araimasu} — \textquotedblleft{}Tomorrow I will wash my face first of all.\textquotedblright{}
\end{quote}

\begin{cousinweb}
\textbf{\ja{まず}} (\emph{mazu}) is a native Japanese word (wago), inherited rather than borrowed from Chinese, and its older origin is not settled. It is usually written in kana; its kanji is a later kun'yomi (JMdict, KANJIDIC2).
\end{cousinweb}

\begin{grammarlens}[title={What you've built}]
One more word for everyday life at A2.
\end{grammarlens}

\subsection*{Your turn}
\begin{itemize}
\raggedright
  \item \emph{Say it:} \emph{mazu}
  \item \emph{Say it:} \emph{mazu}, once more
  \item \emph{Say it:} \emph{ashita wa mazu kao o araimasu} — Tomorrow I will wash my face first of all.
  \item \emph{Recall:} say the Japanese for to spend time, then the Japanese for the beginning, then say \emph{mazu} again
\end{itemize}

\begin{tcolorbox}[breakable,colback=teal!4,colframe=teal!35!black,title={Before you move on}]
What does \ja{まず} mean? (\textquotedblleft{}First of all; to begin with\textquotedblright{}.) Say it once more.
\end{tcolorbox}
\section[\texorpdfstring{\ja{つぎ} (tsugi)}{つぎ (tsugi)}]{\ja{つぎ} (tsugi) — next; the next one}

\label{lesson:JA-C171-tsugi}



\begin{quote}
Before the new one: say the Japanese for the beginning, then the Japanese for first of all.
\end{quote}

\subsection*{You'll want to know: \ja{つぎ}}
\textbf{\ja{つぎ}} — \emph{tsugi} — \textquotedblleft{}next; the next one\textquotedblright{}.

\begin{quote}
\textbf{\ja{つぎは} \ja{らいげつでしょう。}} — \emph{tsugi wa raigetsu deshō} — \textquotedblleft{}The next one will probably be next month.\textquotedblright{}
\end{quote}

\begin{cousinweb}
\textbf{\ja{つぎ}} (\emph{tsugi}) is a native Japanese word (wago), inherited rather than borrowed from Chinese; its kanji was attached later as a kun'yomi, and its older origin is not settled (JMdict, KANJIDIC2).
\end{cousinweb}

\begin{grammarlens}[title={What you've built}]
One more word for everyday life at A2.
\end{grammarlens}

\subsection*{Your turn}
\begin{itemize}
\raggedright
  \item \emph{Say it:} \emph{tsugi}
  \item \emph{Say it:} \emph{tsugi}, once more
  \item \emph{Say it:} \emph{tsugi wa raigetsu deshō} — The next one will probably be next month.
  \item \emph{Recall:} say the Japanese for the beginning, then the Japanese for first of all, then say \emph{tsugi} again
\end{itemize}

\begin{tcolorbox}[breakable,colback=teal!4,colframe=teal!35!black,title={Before you move on}]
What does \ja{つぎ} mean? (\textquotedblleft{}Next; the next one\textquotedblright{}.) Say it once more.
\end{tcolorbox}
\section[\texorpdfstring{\ja{さいご} (saigo)}{さいご (saigo)}]{\ja{さいご} (saigo) — the last; the very end}

\label{lesson:JA-C171-saigo}



\begin{quote}
Before the new one: say the Japanese for first of all, then the Japanese for next.
\end{quote}

\subsection*{You'll want to know: \ja{さいご}}
\textbf{\ja{さいご}} — \emph{saigo} — \textquotedblleft{}the last; the very end\textquotedblright{}.

\begin{quote}
\textbf{\ja{あしたは} \ja{さいごに} \ja{おちゃを} \ja{のみます。}} — \emph{ashita wa saigo ni ocha o nomimasu} — \textquotedblleft{}Tomorrow I will drink tea last.\textquotedblright{}
\end{quote}

\begin{cousinweb}
\textbf{\ja{さいご}} (\emph{saigo}) is Sino-Japanese (kango): two Chinese characters read with their on'yomi, readings borrowed from Chinese. They are \emph{sai}, \textquotedblleft{}most, utmost\textquotedblright{} (Middle Chinese \emph{tswajH}) and \emph{go}, \textquotedblleft{}behind, after\textquotedblright{} (\emph{huwX}) (KANJIDIC2; Middle Chinese in Baxter's transcription).
\end{cousinweb}

\begin{grammarlens}[title={What you've built}]
One more word for everyday life at A2.
\end{grammarlens}

\subsection*{Your turn}
\begin{itemize}
\raggedright
  \item \emph{Say it:} \emph{saigo}
  \item \emph{Say it:} \emph{saigo}, once more
  \item \emph{Say it:} \emph{ashita wa saigo ni ocha o nomimasu} — Tomorrow I will drink tea last.
  \item \emph{Recall:} say the Japanese for first of all, then the Japanese for next, then say \emph{saigo} again
\end{itemize}

\begin{tcolorbox}[breakable,colback=teal!4,colframe=teal!35!black,title={Before you move on}]
What does \ja{さいご} mean? (\textquotedblleft{}The last; the very end\textquotedblright{}.) Say it once more.
\end{tcolorbox}
\section[\texorpdfstring{\ja{そろそろ} (sorosoro)}{そろそろ (sorosoro)}]{\ja{そろそろ} (sorosoro) — soon; any time now}

\label{lesson:JA-C171-sorosoro}



\begin{quote}
Before the new one: say the Japanese for next, then the Japanese for the last.
\end{quote}

\subsection*{You'll want to know: \ja{そろそろ}}
\textbf{\ja{そろそろ}} — \emph{sorosoro} — \textquotedblleft{}soon; any time now\textquotedblright{}.

\begin{quote}
\textbf{\ja{そろそろ} \ja{かえります。}} — \emph{sorosoro kaerimasu} — \textquotedblleft{}I will be going home soon.\textquotedblright{}
\end{quote}

\begin{cousinweb}
\textbf{\ja{そろそろ}} (\emph{sorosoro}) is a native Japanese mimetic word, inherited rather than borrowed from Chinese, and its older origin is not settled (JMdict).
\end{cousinweb}

\begin{grammarlens}[title={What you've built}]
One more word for everyday life at A2.
\end{grammarlens}

\subsection*{Your turn}
\begin{itemize}
\raggedright
  \item \emph{Say it:} \emph{sorosoro}
  \item \emph{Say it:} \emph{sorosoro}, once more
  \item \emph{Say it:} \emph{sorosoro kaerimasu} — I will be going home soon.
  \item \emph{Recall:} say the Japanese for next, then the Japanese for the last, then say \emph{sorosoro} again
\end{itemize}

\begin{tcolorbox}[breakable,colback=teal!4,colframe=teal!35!black,title={Before you move on}]
What does \ja{そろそろ} mean? (\textquotedblleft{}Soon; any time now\textquotedblright{}.) Say it once more.
\end{tcolorbox}
\section[(dialogue)]{Words for days, dates and plans: a review}

\label{lesson:JA-R171-review-170-171}



\begin{quote}
Say the Japanese for the week after next and the year after next.
\end{quote}

\subsection*{Your turn}
Say these aloud:

\begin{itemize}
\raggedright
  \item the week after next, the year after next, the New Year, to celebrate, to spend time
  \item the beginning, first of all, next, the last, soon
\end{itemize}

\begin{tcolorbox}[breakable,colback=teal!4,colframe=teal!35!black,title={Before you move on}]
The week after next? (\textbf{\ja{さらいしゅう}}, \emph{saraishū}.) The year after next? (\textbf{\ja{さらいねん}}, \emph{sarainen}.) The New Year? (\textbf{\ja{しょうがつ}}, \emph{shōgatsu}.) To celebrate? (\textbf{\ja{いわう}}, \emph{iwau}.) To spend time? (\textbf{\ja{すごす}}, \emph{sugosu}.) The beginning? (\textbf{\ja{さいしょ}}, \emph{saisho}.) First of all? (\textbf{\ja{まず}}, \emph{mazu}.) Next? (\textbf{\ja{つぎ}}, \emph{tsugi}.) The last? (\textbf{\ja{さいご}}, \emph{saigo}.) Soon? (\textbf{\ja{そろそろ}}, \emph{sorosoro}.)
\end{tcolorbox}
\section[(dialogue)]{First of All, Next and Last — once more}

\label{lesson:JA-R171-review-final-171}



\begin{quote}
Say the Japanese for the beginning and first of all.
\end{quote}

\subsection*{Your turn}
\begin{itemize}
\raggedright
  \item \emph{Say it:} the beginning, first of all, next, the last, soon
\end{itemize}

\begin{tcolorbox}[breakable,colback=teal!4,colframe=teal!35!black,title={Before you move on}]
The beginning? (\textbf{\ja{さいしょ}}, \emph{saisho}.) First of all? (\textbf{\ja{まず}}, \emph{mazu}.) Next? (\textbf{\ja{つぎ}}, \emph{tsugi}.) The last? (\textbf{\ja{さいご}}, \emph{saigo}.) Soon? (\textbf{\ja{そろそろ}}, \emph{sorosoro}.)
\end{tcolorbox}
